%%%%%%%%%%%%%%%%%%%%%%%%%%%%%%%%%%%%%%%%%%%%%%%%%%%%%%%%%%%%%%%%%%%%%%%%%%
%% Supplementary Material
%% Causal Transfer Learning: Warm-Starting NOTEARS, DAGMA, and GOLEM
%% with Cross-Dataset Graphs  (JBCB-1496)
%%%%%%%%%%%%%%%%%%%%%%%%%%%%%%%%%%%%%%%%%%%%%%%%%%%%%%%%%%%%%%%%%%%%%%%%%%
\documentclass[11pt]{article}

\usepackage[margin=1in]{geometry}
\usepackage{amsmath,amssymb}
\usepackage{graphicx,booktabs,longtable}
\usepackage[super]{cite}
\usepackage[hidelinks]{hyperref}

\newcommand{\to}{$\rightarrow$}
\renewcommand{\thetable}{S\arabic{table}}

\title{\vspace{-2.0em}Supplementary Material\\[4pt]
\large Causal Transfer Learning: Warm-Starting NOTEARS, DAGMA, and GOLEM\\ with Cross-Dataset Graphs}
\author{Shuaidong Gao}
\date{}

\begin{document}
\maketitle
\vspace{-2.5em}

\noindent This document contains supporting analyses, full implementation details, and
per-configuration results that accompany the main text. Section~S3 provides the
edge-stratification table requested in review; Section~S5 reports per-seed
genome-scale results. References to sections (S1--S11) and tables (S1--S12) refer to
this document unless stated otherwise.

\section*{S1. Biological universality of causal programs}

The transfer results rest on a falsifiable hypothesis: causal gene regulatory programs
contain a conserved backbone that recurs across tissue contexts. The hypothesis draws on
three lines of prior evidence.

First, the TCGA Pan-Cancer Atlas established that molecular subtypes---and the
transcriptional programs that define them---transcend tissue-of-origin
boundaries.\cite{hoadley2018cell} If expression-based molecular subtypes are shared, the
regulatory programs that generate those expression patterns may also be shared.

Second, essentiality screens\cite{tsherniak2017} identify a core set of $\sim$1,500
fitness genes whose knockout impairs viability across hundreds of cell lines. The
regulatory logic governing these essential genes is under strong evolutionary constraint
and is unlikely to vary across contexts.

Third, developmental signaling pathways (Wnt, Notch, Hedgehog, TGF-$\beta$) are co-opted
in cancer with conserved downstream transcriptional targets.\cite{zhan2017wnt} The
transferred edges---including WNT2\to C2orf40---belong to precisely these pathways.

Taken together, these observations predict that a subset of regulatory edges should
recur across cancer types, and that this subset should be enriched for core programs
rather than for tissue-specific ones. The warm-start advantage reported in the main text
is consistent with this prediction: the transferred edges are concentrated in conserved
programs, whereas the edges that fail to transfer are predominantly dataset-specific or
weakly supported.

\section*{S2. High-dimensional extension ($d = 5{,}000$): implementation details}

The main text reports that warm-starting preserves $81.8\pm0.4\%$ of source edges at
$d=5{,}000$, versus $45.4\pm2.9\%$ from scratch across three seeds, both evaluated against
each seed's own source graph. This section documents the computational and numerical
choices behind that result.

\paragraph{Optimizer.} At $d=5{,}000$ the parameter matrix has $25$ million entries whose
per-parameter curvature history would be memory-prohibitive for L-BFGS-B, so the inner
optimizer is ADAM ($\mathrm{lr}=0.002$, full-batch gradients) rather than the L-BFGS-B
used at $d\le889$. The augmented-Lagrangian outer loop uses 20 iterations with
$\rho_0=1$, a larger initial penalty than the $\rho_0=0.1$ of the main experiments, chosen
to accelerate DAG enforcement at this scale. All other hyperparameters ($\lambda=0.01$,
edge threshold $\tau=0.3$) match the main experiments.

\paragraph{Numerical stability.} The acyclicity term
$\operatorname{tr}(\exp(W\circ W))$ is evaluated in float32, where the matrix exponential
overflows once the row-sum of $W\circ W$ approaches $\sim$30. Source-graph weights are
therefore bounded so that this row-sum stays below that bound (here $1.39$), and scratch
training uses random initialization $\mathcal{N}(0,0.05^2)$ across seeds (42, 123, 456),
reduced from the $0.1$ used at $d\le889$ because $0.1$ would place the row-sum near $50$.
Each seed learned its own source graph from that random initialization, giving
$624\pm4$ source edges at $\tau=0.3$ (range 619--627); the warm-start run for a seed
uses that seed's own source graph, as in the $d\le889$ benchmark. For reference, a single
deterministic run initialized at $W=0$ yields a denser 634-edge source graph and a warm
retention of $513/634=80.9\%$, consistent with the seed-matched estimate.

\paragraph{Compute.} NOTEARS converged on a single consumer GPU (NVIDIA RTX 5060 Laptop,
8~GB) using $7.7$~GB of VRAM and $3.5$~hours per run, with the $\mathcal{O}(d^3)$ matrix
exponential as the sole computational bottleneck. Warm-starting also recovers more total
edges than scratch ($842\pm10$ vs.\ $703\pm9$).

\paragraph{Relation to the scaling law.} Retention is necessarily bounded by $100\%$, so
the power law $r(d)\propto d^{0.47}$ established at $d\le889$ cannot be extrapolated
indefinitely; the measured $d=5{,}000$ warm retention ($81.8\pm0.4\%$) is accordingly below
a naive extrapolation, while the absolute advantage over scratch ($+36.4$ percentage
points, versus $+34.1$~pp at $d=889$) remains substantial and consistent with the trend of
growing transfer benefit with dimension.

\paragraph{Extrapolation beyond $d=5{,}000$.} Extrapolating the cubic time scaling from
the measured run gives $d=10{,}000$ at $\sim$28~hours ($8\times$ the $d=5{,}000$ runtime)
and $d=16{,}000$ at roughly 5~days. Assuming quadratic memory scaling---a conservative
upper bound---$d=10{,}000$ would require $\sim$31~GB and $d=16{,}000$ $\sim$79~GB, beyond
an 8-GB consumer GPU but within reach of 24-GB or multi-GPU systems. Low-rank NOTEARS with
rank $r\approx100$ reduces the bottleneck to $\mathcal{O}(d\,r^2)$, bringing genome-scale
causal transfer within reach of a single GPU.

\paragraph{Why the log-determinant term is rigid.} The main text compresses the
DAGMA mechanism to three sentences; the full argument is as follows. DAGMA's
characteristic loss term $-\log|\det(I-W)|$ contributes the gradient
$\nabla_W[-\log|\det(I-W)|]=(I-W)^{-\top}$, which is bounded at a feasible source
solution: for a DAG the matrix $I-W_S$ is unit-triangular in topological order, so
every eigenvalue equals one and $\det(I-W_S)=1$. Any perturbation that would close a
cycle drives $\det(I-W)\to 0$, where the inverse---and hence the gradient---diverges.
The log-det term therefore acts as a one-sided barrier: remaining at $W_S$ is cheap,
leaving it is expensive, which is precisely the local rigidity that yields
dimension-independent retention. (Acyclicity itself is enforced by the
trace-exponential constraint; the log-det term is the additional penalty that shapes
the local geometry.) The MSE loss of NOTEARS has the opposite geometry: its Hessian,
$\tfrac{1}{n}X^\top X$ up to the DAG penalty, has a spectrum spanning orders of
magnitude on transcriptomic covariates, so the objective is nearly flat along many
directions and admits a combinatorially growing set of near-equivalent solutions as
$d$ increases. Warm-starting thus supplies a bias that the MSE objective cannot
overturn, whereas the augmented Lagrangian term $h(W)^2$ exerts only gradual pressure
toward the DAG constraint. The power-law exponent $0.47$ follows: as $d$ grows, the
number of near-equivalent MSE solutions increases combinatorially, so the warm-start
bias becomes proportionally more valuable.

\paragraph{Pathway mapping of the seven transferred DepMap edges.} The DepMap$\to$BRCA
transfer retains seven edges at $\tau=0.3$ (35\% of the 20 source edges; full list in
\texttt{results/depmap\_edges\_d454.json}). Their gene annotations map to calcium
homeostasis (S100P$\to$PPP1R14C), epithelial differentiation (POF1B$\to$KRT16), cell
adhesion and metastasis (SLC44A4$\to$TSPAN8$\to$PCSK9), Notch signaling and
antioxidant defence (HES2$\to$GPX2), protein processing (PCSK9$\to$SERPINB5) and
ubiquitin-mediated proteolysis (USP9Y$\to$POF1B). The calcium-homeostasis edge
S100P$\to$PPP1R14C is preserved under both the TCGA and DepMap transfers, marking a
core relationship conserved across measurement platforms.

\section*{S3. Edge stratification and pathway membership}

Table~S3 stratifies the 73 source edges of the TCGA$\to$BRCA transfer at
$\tau=0.3$ by source-graph weight $|W_S|$. For each band it reports the number of
edges, the retention achieved by scratch and by warm-start at the same threshold,
and the KEGG pathways most represented among the genes touched by the band.
Warm-start matches or exceeds scratch in every band. KEGG annotation was resolved
through the KEGG REST API (\texttt{/find/hsa/} for symbol$\to$entrez and
\texttt{/link/pathway/} for pathway membership); 51 of the 105 genes touched by
source edges carried at least one KEGG pathway annotation.

\setcounter{table}{2}
\begin{table}[h]
\small
\caption{Edge-strength stratification at $\tau=0.3$ ($d=454$, TCGA$\to$BRCA; 73
source edges). Retention is the fraction of band edges recovered at $|W|>0.3$.}
\begin{tabular}{@{}lcccc@{}}
\toprule
Band $|W_S|$ & $n$ & Scratch (\%) & Warm (\%) & $\Delta$ (pp)\\
\midrule
$[0.3,0.4)$ & 25 & 24.0 & 44.0 & $+20.0$\\
$[0.4,0.5)$ & 18 & 16.7 & 33.3 & $+16.6$\\
$[0.5,0.6)$ & 8 & 62.5 & 62.5 & $0.0$\\
$[0.6,0.7)$ & 10 & 80.0 & 90.0 & $+10.0$\\
$[0.7,1.0)$ & 12 & 50.0 & 83.3 & $+33.3$\\
\midrule
All & 73 & 38.4 & 56.2 & $+17.8$\\
\bottomrule
\end{tabular}
\label{tab:S3}
\end{table}

\noindent\textbf{Predominant KEGG pathways per band (gene counts in parentheses).}
Full per-gene annotation is stored in \texttt{results/edge\_stratification\_tau03.json}.
\begin{itemize}\itemsep2pt
\item $[0.3,0.4)$: Metabolic pathways (6), Cornified envelope formation (4),
Cadherin signaling (4), Chemical carcinogenesis--reactive oxygen species (3),
Protein digestion and absorption (2), Glutathione metabolism (2), Estrogen
signaling (2), \textit{Staphylococcus aureus} infection (2).
\item $[0.4,0.5)$: Metabolic pathways (4), Cornified envelope formation (3),
Retinol metabolism (2), Metabolism of xenobiotics by cytochrome P450 (2), Drug
metabolism--cytochrome P450 (2), Cytokine--cytokine receptor interaction (2),
Prostate cancer (2), Cadherin signaling (2).
\item $[0.5,0.6)$: Cornified envelope formation (2), Cadherin signaling (2),
Estrogen signaling (2), \textit{Staphylococcus aureus} infection (2),
TGF-$\beta$ signaling (1), Glycosaminoglycan biosynthesis--heparan sulfate (1),
Renin--angiotensin system (1), Endocrine calcium reabsorption (1).
\item $[0.6,0.7)$: Cornified envelope formation (3), Cadherin signaling (3),
Transcriptional misregulation in cancer (2), MicroRNAs in cancer (2), Estrogen
signaling (2), \textit{Staphylococcus aureus} infection (2), Choline metabolism in
cancer (2), Cytoskeleton in muscle cells (1).
\item $[0.7,1.0)$: Metabolic pathways (3), Cadherin signaling (3), Mucin type
O-glycan biosynthesis (1), Glycosphingolipid biosynthesis (1), Lysine degradation
(1), Neuroactive ligand--receptor interaction (1), Neuroactive ligand signaling
(1), Retrograde endocannabinoid signaling (1).
\end{itemize}

Pathway membership is reported as annotation rather than as an enrichment
statistic: only two of the 73 source edges touch the canonical \emph{Pathways in
cancer} map (hsa05200), so this stratification is not powered for pathway-level
claims, and the main text accordingly reports transfer rates by edge strength
rather than by pathway.

\section*{S4. External validation: database-level detail}

Validation was performed on the top-200 edges by absolute weight for each method
($d=454$). STRING v12 was queried at combined-score thresholds 400, 700 and 900; mapping
succeeded for $97.8\%$ of genes. TRRUST, RegNetwork and PathwayCommons are curated
resources recording experimentally supported transcription-factor/target and
protein--protein interactions, and therefore overlap only marginally with
co-expression-inferred associations.

\setcounter{table}{3}
\begin{table}[h]
\caption{External validation detail ($d=454$, top-200 edges per method). Percentages
are the fraction of mappable edges supported by each database.}
\begin{tabular}{@{}lccccc@{}}
\toprule
Method & STRING $\ge400$ & STRING $\ge700$ & STRING $\ge900$ & BioGRID & TRRUST\\
\midrule
TCGA source & 4.6 & 3.0 & 2.0 & 0.5 & 0.0\\
Scratch     & 4.5 & 3.0 & 2.5 & 1.0 & 0.0\\
TCGA-warm   & 4.5 & 3.0 & 2.5 & 1.0 & 0.0\\
DM-warm     & 4.5 & 3.0 & 2.5 & 1.0 & 0.0\\
\bottomrule
\end{tabular}
\label{tab:Sstring}
\end{table}

BioGRID v5.0.260 was queried for the same edge sets, yielding $1.0\%$ precision
(2/200 edges) for scratch, TCGA-warm and DepMap-warm, and $0.5\%$ (1/200) for the TCGA
source graph. The two BioGRID-confirmed pairs are KRT16--KRT7 and PRSS3--AGR2. RegNetwork
and PathwayCommons yield $0\%$ across all methods, as does TRRUST. The five preserved
edges central to the cross-domain analysis (WNT2\to C2orf40, S100P\to PPP1R14C,
PRKY\to ATP2C2, LGI3\to SLC26A9, POF1B\to SLC44A4) are absent from BioGRID, consistent
with their status as \emph{de novo} discoveries.

The relevant comparison is not the absolute precision but its invariance across methods:
scratch and both transfer variants reach essentially the same precision at every
threshold. Had warm-starting manufactured spurious edges, precision would have fallen;
had it merely re-derived database-annotated edges, precision would have risen. Neither
occurs.

\section*{S5. Per-seed genome-scale results}

Table~S5 reports the $d=5{,}000$ experiment at the level of individual seeds. The
protocol matches the $d\le889$ benchmark exactly: each seed (i) learns its own source
NOTEARS graph on the 9,241 non-breast samples from a random initialization
$\mathcal{N}(0,0.05^2)$, and (ii) runs both a warm-start (initialized at that source
graph) and a scratch run (initialized from the same distribution) on the 1,218 BRCA
samples. Retention is the fraction of that seed's source edges recovered at
$|W|>0.3$, so numerator and denominator are always seed-matched; the earlier single-run
figure of $513/634=80.9\%$ is superseded by this protocol. Because the source graph
itself differs across seeds ($619$--$627$ edges), the denominator is not fixed, which is
the same reason the $d\le889$ tables carry non-zero variance in both columns.

\setcounter{table}{4}
\begin{table}[h]
\small
\caption{Per-seed genome-scale transfer ($d=5{,}000$, $\tau=0.3$). Retention for every
method is measured against the same seed's source graph.}
\begin{tabular}{@{}lcccccc@{}}
\toprule
Seed & Source edges & Warm edges & Warm kept (ret.) & Scratch edges & Scratch kept (ret.) & $\Delta$ (pp)\\
\midrule
42  & 619 & 856 & 505 (81.6\%) & 702 & 299 (48.3\%) & $+33.3$\\
123 & 626 & 833 & 510 (81.5\%) & 693 & 290 (46.3\%) & $+35.1$\\
456 & 627 & 838 & 516 (82.3\%) & 714 & 260 (41.5\%) & $+40.8$\\
\midrule
Mean$\pm$SD & $624\pm4$ & $842\pm10$ & $81.8\pm0.4$\% & $703\pm9$ & $45.4\pm2.9$\% & $+36.4$\\
\bottomrule
\end{tabular}
\label{tab:S5genome}
\end{table}

Two observations follow. First, warm retention is remarkably stable across seeds
(SD $0.4$~pp), whereas scratch retention varies by $6.8$~pp: the warm-start solution is
anchored by the source structure, while scratch converges to a different near-equivalent
DAG on each run. Second, the per-seed gain is positive for every seed and increases as
scratch degrades ($+33.3$, $+35.1$, $+40.8$~pp), i.e.\ the benefit is largest precisely
when the unguided optimization lands furthest from the source structure.

\paragraph{Relation to the single-run figures in the first response.} The first response
reported a single warm-start value, $513/634=80.9\%$, and a three-seed scratch value of
$38.1\pm1.4\%$. Both were measured against one deterministic source graph obtained from a
zero initialization; the warm-start figure consequently had no seed-to-seed variance, and
neither figure corresponds to the seed-matched denominator used throughout the
$d\le889$ tables. Evaluating the present warm runs against that same 634-edge reference
instead gives $65.3\pm0.7\%$ versus $38.1\pm1.4\%$ (+27.2~pp), lower than the seed-matched
$81.8\%$ because each seed's own source graph overlaps the reference graph by only
77--79\%, and $0.78\times0.82\approx0.65$. The seed-matched figures in Table~S5 are the
ones that answer the reviewer's request, since they are the quantities reported for every
other dimension; the direction and the significance of the conclusion are unchanged under
either denominator, all three seeds gaining 27--41 percentage points. Per-seed matrices
and the consolidated statistics are stored in \texttt{results/d5000\_seedmatched.json},
\texttt{results/d5000\_multiseed\_transfer.json} and
\texttt{results/d5000\_checkpoints/}.

\section*{S6. Reproducibility}

All experiments use fixed random seeds (42, 123, 456). Data are public: TCGA via the
UCSC Xena browser (\url{https://xenabrowser.net/}), DepMap
(\url{https://depmap.org/}), STRING v12 (\url{https://string-db.org/}), and TRRUST v2
(\url{https://www.grnpedia.org/trrust/}). Source code and pre-computed causal graphs are
available at \url{https://github.com/sgao-academics/causal-transfer-discovery}.
The main-text analyses are reproduced by \texttt{scripts/run\_all.py}; the genome-scale
extension by \texttt{scripts/run\_d5000\_multiseed.py}.

\section*{S7. Sensitivity to the edge threshold}

\setcounter{table}{6}
\begin{table}[h]
\caption{Edge retention at varying thresholds ($d=454$, TCGA$\to$BRCA). Counts are
source-graph edges at that threshold that are also present in the scratch or warm
graph at the same threshold.}
\begin{tabular}{@{}lcccc@{}}
\toprule
$\tau$ & Source edges & Scratch overlap & Warm overlap & $\Delta$\\
\midrule
0.1 & 1156 & 122 (11\%) & 158 (14\%) & $+36$\\
0.2 & 168 & 38 (23\%) & 57 (34\%) & $+19$\\
0.3 & 73 & 28 (38\%) & 41 (56\%) & $+13$\\
0.4 & 48 & 20 (42\%) & 28 (58\%) & $+8$\\
0.5 & 30 & 13 (43\%) & 20 (67\%) & $+7$\\
\bottomrule
\end{tabular}
\label{tab:S7}
\end{table}

Pooling the $\tau\ge0.5$ bands of Table~S3 (the 30 source edges with $|W_S|>0.5$)
reproduces the $\tau=0.5$ row of Table~S7 (warm $20/30=67\%$, scratch
$13/30=43\%$), confirming that the strength-stratified and threshold-stratified
views are consistent.

\section*{S8. Paired statistics for the core comparison}

Every main-text comparison is paired: the warm-start and scratch runs for a given
$(d,\text{seed})$ share the same gene set, the same target samples and the same
optimizer, so the appropriate test is a paired one. Table~S8 gives the per-dimension
paired differences for NOTEARS. Across the 15 NOTEARS pairs the mean gain is $+26.3$~pp
(Wilcoxon signed-rank $p=6.1\times10^{-5}$; exact sign test 15/15, $p=3.1\times10^{-5}$),
and the smallest single gain is $+7.9$~pp at $d=200$.

\paragraph{DAGMA and GOLEM use the same protocol.} The two additional algorithms in
Table~4.4 of the main text are run identically: each seed $s\in\{42,123,456\}$ learns its
own source graph on the 9,241 non-breast samples, the warm arm is initialised from that
graph, the scratch arm is initialised from $\mathcal{N}(0,0.1^2)$, and both are scored
against the \emph{same} source graph. The denominators are therefore seed-matched, as for
NOTEARS; the three seeds differ in the source graph and hence carry genuine variance in
both arms. This is what makes the per-dimension standard deviations in Table~4.4
meaningful and what makes the 45 configurations 45 independent comparisons. Across all 45
algorithm--dimension--seed configurations warm-start exceeds scratch (sign test
45/45, $p=2.8\times10^{-14}$). The shuffled control, by contrast, is indistinguishable
from scratch (paired Wilcoxon $p=0.86$, mean difference $+0.2$~pp), while warm-start
exceeds shuffled at all 15 pairs ($p=6.1\times10^{-5}$).

\setcounter{table}{7}
\begin{table}[h]
\small
\caption{Per-dimension paired differences (NOTEARS, $\tau=0.3$, 3 seeds). The
difference column is warm minus scratch within the same $(d,\text{seed})$.}
\begin{tabular}{@{}lcccccc@{}}
\toprule
$d$ & Scratch (\%) & Warm (\%) & Paired $\Delta$ (pp) & min & max & wins\\
\midrule
200 & 19.9 & 31.9 & $+12.0$ & $+7.9$ & $+17.2$ & 3/3\\
454 & 33.4 & 55.1 & $+21.7$ & $+17.8$ & $+24.4$ & 3/3\\
600 & 33.3 & 66.6 & $+33.3$ & $+27.6$ & $+39.3$ & 3/3\\
750 & 36.5 & 66.8 & $+30.3$ & $+25.7$ & $+35.6$ & 3/3\\
889 & 34.0 & 68.1 & $+34.0$ & $+29.9$ & $+36.5$ & 3/3\\
\midrule
All & --- & --- & $\mathbf{+26.3}$ & $+7.9$ & $+39.3$ & $\mathbf{15/15}$\\
\bottomrule
\end{tabular}
\label{tab:S8paired}
\end{table}

\paragraph{Standard-deviation convention.} All standard deviations in the main text and
in this document are sample standard deviations ($n-1$). Reporting the population
convention ($n$) would reduce every value by the factor
$\sqrt{(n-1)/n}=0.816$ (18\% at $n=3$); we state the convention explicitly, and every
mean, SD and paired difference can be recomputed from
\texttt{results/benchmark\_v6\_checkpoint.json}.

\paragraph{Post-hoc thresholding.} The acyclicity constraint is enforced on the
unthresholded matrix, where $h(W)<10^{-6}$ by construction. Thresholding at $\tau=0.3$
is applied only to obtain an interpretable edge set for edge counts, retention and
database validation, and (as is generally the case for post-hoc thresholding of
continuous DAG estimators) the thresholded graph need not itself be acyclic.


\section*{S9. Is pooling a substitute for transfer?}

If the source samples are available, a referee will ask why the cohorts are not simply
concatenated and one graph fitted to the union. We ran that arm: the 9{,}241 source and
1{,}218 BRCA samples are stacked, standardised jointly, and a single NOTEARS graph is
estimated with the same penalty, threshold and outer-loop budget as every other arm, then
scored against the same source graph.

\paragraph{Retention is not comparable across arms here.} Pooling sees 88\% of the samples
that generated the source graph, so its retention is high by construction; in the other
direction, $\tau=0.3$ is an absolute cut on weights whose scale shrinks as $n$ grows, so the
number of edges surviving the cut falls as samples are added. Neither effect reflects
structural agreement. The comparison that survives both problems is the held-out residual of
the thresholded graph on the 20\% of BRCA withheld from every arm, computed after every graph
has been cut to the \emph{same} number of edges $k$: at fixed $k$ the residual measures how
well each edge set explains breast-tissue covariance, and the complexity penalty is identical
by construction.

\setcounter{table}{8}
\begin{table}[h]
\caption{Pooling the source and target samples is not a substitute for transfer. One graph is fitted to the stacked matrix of all 9{,}241 source and 1{,}218 BRCA samples and scored against the same source graph. Retention is not comparable across arms here, because $\tau=0.3$ is an absolute cut on weights whose scale depends on the sample size; the fair comparison is the held-out residual at a matched edge count $k$ (last two columns, lower is better). Seed 42 shown; all dimensions share that seed.}
\begin{tabular}{@{}lrrrrr@{}}
\toprule
$d$ & Source edges & Warm ret.\ (\%) & Pool ret.\ (\%) & Warm NRMSE & Pool NRMSE\\
\midrule
200 & 64 & 29.7 & 54.7 & 0.8578 & 0.8850\\
454 & 74 & 56.8 & 56.8 & 0.9251 & 0.9328\\
600 & 85 & 64.7 & 51.8 & 0.9386 & 0.9419\\
750 & 101 & 69.3 & 52.5 & 0.9498 & 0.9533\\
889 & 107 & 67.3 & 56.1 & 0.9540 & 0.9579\\
\bottomrule
\end{tabular}
\label{tab:S9}
\end{table}

\noindent Across all 11 algorithm--dimension--seed cells the pooled graph has a higher held-out residual than the warm graph by $+0.012$ on average (range $+0.003$ to $+0.031$); the penalty is positive in 11 of 11 cells.

\paragraph{Target-only priors do not reproduce the gain.} Three further arms build a starting
matrix without the source cohort: the $|\text{corr}|$ skeleton of BRCA itself, its ridge
coefficients, and a random DAG matching the source graph's edge count and out-degree profile.
None approaches warm-start's retention, and the degree-matched control tracks scratch almost
exactly---so the benefit cannot be attributed to ``any sparse hint''.

\setcounter{table}{9}
\begin{table}[h]
\caption{Target-only structure priors do not reproduce the transfer benefit. All arms start NOTEARS from a matrix built without the source cohort and are scored against the source graph; means over three seeds.}
\begin{tabular}{@{}lrrrrrrrr@{}}
\toprule
$d$ & Scratch & Warm & Shuffled & Pool & Corr-prior & Ridge-prior & Degree-matched & $k$\\
\midrule
200 & 19.9 & 31.9 & 17.9 & 56.0 & 17.3 & 16.8 & 21.5 & 66\\
454 & 33.4 & 55.1 & 35.6 & 55.5 & 33.8 & 33.7 & 36.0 & 94\\
\bottomrule
\end{tabular}
\label{tab:S10}
\end{table}

\noindent At $k=70$ the six source-free or sparse priors differ from one another by at most $0.007$ in held-out residual within a cell (median spread $0.004$), whereas the pooled graph sits $+0.008$ above their mean (median). Predictive residual therefore does not discriminate among sparse priors, but it does penalise the pooled one.

\section*{S10. Dependence on the size of the target cohort}

The motivating claim of the paper is that transfer matters most when the target cohort is
small. Table~S11 tests it by subsampling BRCA to $n=150$, $300$, $600$ and the full
$1{,}218$ samples, holding the source graph, the optimiser, the penalty and the threshold
fixed. Warm-start retention is almost flat across a factor of eight in $n$
($55$--$58\%$), whereas scratch rises from $21\%$ to $34\%$: the advantage is $+37.2$~pp at
$n=150$ against $+21.2$~pp at $n=1{,}218$, and it is positive for every seed at every $n$.
What transfer supplies is therefore not extra samples but a starting point whose quality does
not depend on how many target samples happen to be available.

\setcounter{table}{10}
\begin{table}[h]
\caption{Retention versus target sample size ($d=454$, TCGA$\to$BRCA). Each $n$ draws a random subset of the 1{,}218 BRCA samples; the source graph, optimiser and threshold are unchanged. Mean $\pm$ sample SD over three seeds.}
\begin{tabular}{@{}rrrr@{}}
\toprule
$n$ & Scratch (\%) & Warm (\%) & $\Delta$ (pp)\\
\midrule
150 & 21.0$\pm$3.4 & 58.2$\pm$4.8 & $+37.2$\\
300 & 33.0$\pm$6.0 & 56.0$\pm$1.7 & $+23.0$\\
600 & 32.1$\pm$2.3 & 57.8$\pm$3.9 & $+25.8$\\
1218 & 33.9$\pm$5.1 & 55.1$\pm$1.7 & $+21.2$\\
\bottomrule
\end{tabular}
\label{tab:S11}
\end{table}

\section*{S11. Ten-seed robustness}

Table~1 rests on three seeds. Table~S12 repeats the NOTEARS comparison on those three plus
seven fresh seeds whose random starts are seeded explicitly, leaving the three main-table
seeds untouched so that no number in the main text changes. The sign of the warm--scratch
difference is unanimous in every configuration reported, and the magnitudes ($+6.8$ to
$+29.2$~pp at $d=454$) bracket the three-seed means of the main text.

\setcounter{table}{11}
\begin{table}[h]
\caption{Ten-seed NOTEARS comparison. The first three seeds are those of the main tables; the remaining seven are fresh draws whose source graphs are seeded explicitly. Each row is one seed's own source graph.}
\begin{tabular}{@{}lrrrrrr@{}}
\toprule
& \multicolumn{3}{c}{$d=454$} & \multicolumn{3}{c}{$d=889$}\\
\cmidrule(lr){2-4}\cmidrule(lr){5-7}
Seed & Scratch & Warm & $\Delta$ & Scratch & Warm & $\Delta$\\
\midrule
42 & 32.4 & 56.8 & $+24.4$ & 30.8 & 67.3 & $+36.5$\\
123 & 35.6 & 53.4 & $+17.8$ & 33.9 & 69.6 & $+35.7$\\
456 & 32.1 & 55.1 & $+23.0$ & 37.4 & 67.3 & $+29.9$\\
789 & 40.3 & 56.9 & $+16.6$ & 48.1 & 64.8 & $+16.7$\\
1011 & 46.6 & 53.4 & $+6.8$ & 47.3 & 67.3 & $+20.0$\\
2024 & 40.0 & 54.7 & $+14.7$ & 51.8 & 67.0 & $+15.2$\\
3141 & 48.6 & 58.3 & $+9.7$ & 44.5 & 71.8 & $+27.3$\\
4242 & 31.9 & 61.1 & $+29.2$ & 47.7 & 71.0 & $+23.3$\\
5555 & 48.5 & 61.8 & $+13.3$ & 46.7 & 69.2 & $+22.5$\\
6789 & 38.2 & 56.6 & $+18.4$ & 45.6 & 66.7 & $+21.1$\\
\bottomrule
\end{tabular}
\label{tab:S12}
\end{table}

\noindent At $d=454$ warm-start exceeds scratch in 10 of 10 seeds (exact sign test $p=2^{-10}=9.8\times10^{-4}$); at $d=889$, in 10 of 10. Retention rises with dimension for both arms, reproducing the scaling behaviour of the main text on independent draws.

\bibliographystyle{plain}
\bibliography{refs}

\end{document}
